\subsection*{Determinación de la prolongación analítica correlacionada}

La construcción completa de la carta analítica fija el siguiente orden
de dependencias. El estado enriquecido $\mathsf s$ publica primero la
precoordenada
\[
 a(\mathsf s)=p(\mathsf s)+e_0(\mathsf s)-f(\mathsf s)
              -\kappa_{\rm per}(\mathsf s),
\]
donde $p=\pi_{\rm HMT}-3$, $e_0=e_{\rm HMT}-2$ y
$f=\varphi_{\rm HMT}-1$. Las tres coordenadas regionales y la lectura
periódica del registro pertenecen al mismo desarrollo generativo.
La normalización entera y la carta analítica realizan dos publicaciones
correlacionadas de esta precoordenada.

Con $q=1/729$ y $E_{9,\mathsf s}$ el jet de orden nueve, el coeficiente
de enlace es
\[
 \lambda(\mathsf s)=
 -\frac{E_{9,\mathsf s}(a(\mathsf s))
         \bigl(1-q\,a(\mathsf s)^3\bigr)}{a(\mathsf s)^{10}}.
\]
La prolongación completa queda entonces determinada por
\[
 E_{\mathsf s,\infty}(X)=E_{9,\mathsf s}(X)
      +\lambda(\mathsf s)\frac{X^{10}}{1-X^3/729},
\]
o, equivalentemente, por la recurrencia de todos sus coeficientes:
\[
 b_{10+3n}=729^{-n}\lambda(\mathsf s),\qquad
 b_{11+3n}=b_{12+3n}=0\qquad(n\geq0).
\]

La unicidad de $\lambda$ tiene como dominio la clase de prolongaciones
$E_{9,\mathsf s}(X)+\mu X^{10}/(1-qX^3)$: la condición de anularse en
$a(\mathsf s)$ determina exactamente $\mu=\lambda(\mathsf s)$.
La naturalidad de los truncamientos, el resto geométrico y las cotas de
la derivada establecen después la convergencia y la raíz simple única.
Esta raíz coincide con la publicación entera a toda profundidad.

Éste es el estatuto preciso de la doble realización: identidad de las
publicaciones aritmética y analítica del mismo estado, con dependencia
explícita de la precoordenada generada y de la clase geométrica de
prolongación. La construcción está desarrollada íntegramente en la
sección~\ref{x_reciprocidad:subsec:alpha-lector-completo-rev08},
con sus pruebas de naturalidad, convergencia y compatibilidad.
